\SourceLeaf{355}

\begin{center}
\large\bfseries GEOMETRISCHE SEITE DER TERNÄREN FORMEN.
\end{center}



Ein Punkt $(0)$ im Raum sei als Anfangspunkt angenommen. Die Verschiebungen von dort zu drei
anderen Punkten $\boldsymbol{P}$, $\boldsymbol{P}'$, $\boldsymbol{P}''$, die mit ihm nicht
in einer Ebene liegen, seien der Reihe nach $t$, $t'$, $t''$. Wo keine Verwechslung möglich
ist, können die Punkte $\boldsymbol{P}$, $\boldsymbol{P}'$, $\boldsymbol{P}''$ selbst mit $(t)$,
$(t')$, $(t'')$ bezeichnet werden.

Ferner sei $(t,t')$ allgemein das Produkt der Längen (kürzesten Entfernungen) der beiden
Linien $t$, $t'$ und des Kosinus ihrer Neigung usw.

Man erhält allgemein $(\alpha t + \alpha' t' + \alpha'' t'' + .\,.\,.,\ \beta u + \beta' u' +
\beta'' u'' + .\,.\,.)$, wenn man die Multiplikation
\[
(\alpha t + \alpha' t' + \alpha'' t'' + .\,.\,.) \times (\beta u + \beta' u' + \beta'' u'' + .\,.\,.)
\]
ausführt und $tu$, $tu'$, $tu''$, $t'u'$, $t'u''$ usw. durch $(t,u)$, $(t,u')$,
$(t,u'')$, $(t',u')$, $(t',u'')$ usw. ersetzt.

Jeder Punkt im Raum lässt sich durch ein Trinom
\[
(xt + x't' + x''t'')
\]
darstellen.

Für alle Punkte, die in einer bestimmten Ebene liegen, gilt dann eine Gleichung
\[
\lambda x + \lambda' x' + \lambda'' x'' = L
\]
wobei $\lambda$, $\lambda'$, $\lambda''$, $L$ bestimmte Zahlen bezeichnen. Für eine
Ebene durch die drei Punkte $\mu t$, $\mu' t'$, $\mu'' t''$ gilt
\[
\lambda\mu = \lambda'\mu' = \lambda''\mu'' = L.
\]
